\documentclass[11pt]{article}

\usepackage[utf8]{inputenc}
\usepackage[T1]{fontenc}
\usepackage{lmodern}
\usepackage{geometry}
\usepackage{amsmath,amssymb,amsfonts,amsthm,mathtools}
\usepackage{bm}
\usepackage{enumitem}
\usepackage{microtype}
\usepackage{hyperref}
\usepackage{titlesec}
\usepackage{booktabs}
\usepackage{array}
\usepackage{setspace}
\usepackage{mathrsfs}
\usepackage{physics}

\geometry{margin=1in}
\hypersetup{colorlinks=true, linkcolor=black, urlcolor=blue, citecolor=black}
\setlist[enumerate]{topsep=0.4em,itemsep=0.35em}
\setlist[itemize]{topsep=0.3em,itemsep=0.25em}
\titleformat{\section}{\large\bfseries}{\thesection.}{0.5em}{}
\titleformat{\subsection}{\normalsize\bfseries}{\thesubsection.}{0.5em}{}
\setstretch{1.06}

\newcommand{\Aether}{\AE{}ther}
\newcommand{\Aflow}{\AE{}ther-flow}
\newcommand{\TheoryName}{\Aflow{} Interpretation of Relativity}
\newcommand{\TheoryProgram}{\Aether{} / \Aflow{} framework}

\newcommand{\CompletedMetric}{g^{\sharp}}
\newcommand{\MatterSpace}{\Psi}

\newcommand{\Car}{\mathrm{car}}
\newcommand{\Cur}{\mathrm{cur}}
\newcommand{\Hom}{\mathrm{hom}}

\newcommand{\ForSpace}[1]{\mathcal I_{#1}^{\mathrm{for}}}
\newcommand{\PrimShell}{\mathcal S_{\mathrm{prim}}}
\newcommand{\Reduction}[1]{\mathcal R_{#1}}
\newcommand{\CurrentCarrier}[1]{\mathfrak C_{#1}^{\mathrm{cur}}}
\newcommand{\EqCarrier}[1]{\mathfrak C_{#1}^{\mathrm{eq}}}
\newcommand{\MetricOut}{\mathscr G_{\mathrm{out}}}
\newcommand{\MatterOut}{\mathscr T_{\mathrm{out}}}

\newcommand{\SubTarget}[1]{E_{#1}^{\mathrm{sub}}}
\newcommand{\EqnTarget}[1]{Q_{#1}^{\mathrm{eqn}}}
\newcommand{\HiddenSpace}[1]{\mathcal H_{#1}^{\mathrm{hid}}}
\newcommand{\HiddenBody}[1]{D_{#1}^{\mathrm{hid}}}

\newcommand{\SeedState}[1]{\mathcal S_{#1}^{\mathrm{seed}}}
\newcommand{\SeedBody}[1]{\mathcal B_{#1}^{\mathrm{seed}}}
\newcommand{\SeedShell}[1]{\mathcal Z_{#1}^{\mathrm{seed}}}
\newcommand{\SeedSupportShell}[1]{\mathcal U_{#1}^{\mathrm{seed,sup}}}
\newcommand{\SeedZeroCore}[1]{\mathcal Z_{#1}^{0,\mathrm{seed}}}
\newcommand{\SeedSplit}[1]{\Phi_{#1}^{\mathrm{seed}}}
\newcommand{\SeedCharge}[1]{\mathcal Q_{#1}^{\mathrm{seed}}}
\newcommand{\SeedEmbed}[1]{\zeta_{#1}^{\mathrm{seed}}}
\newcommand{\SeedKernel}[1]{\widetilde K_{#1}^{0,\mathrm{seed}}}
\newcommand{\SeedCoercive}[1]{P_{#1}^{\mathrm{seed,co}}}

\newcommand{\GermNucleus}{\mathfrak G^{\mathrm{sub,germ,zstn}}}
\newcommand{\GermState}[1]{\mathcal G_{#1}^{\mathrm{germ}}}
\newcommand{\GermBody}[1]{\mathcal B_{#1}^{\mathrm{germ}}}
\newcommand{\GermProj}[1]{\Pi_{#1}^{\mathrm{germ}}}
\newcommand{\GermSeedSec}[1]{\sigma_{#1}^{\mathrm{germ\leftarrow seed}}}
\newcommand{\GermSection}[1]{\Gamma_{#1}^{\mathrm{germ}}}
\newcommand{\GermShell}[1]{\mathcal Z_{#1}^{\mathrm{germ}}}

\newcommand{\ProtoTransportCore}{\mathfrak P^{\mathrm{sub,proto,xtc}}}
\newcommand{\ProtoSectionCore}{\mathfrak P^{\mathrm{sub,proto,sec}}}
\newcommand{\ProtoState}[1]{\mathcal P_{#1}^{\mathrm{proto}}}
\newcommand{\ProtoBody}[1]{\mathcal B_{#1}^{\mathrm{proto}}}
\newcommand{\ProtoProj}[1]{\Pi_{#1}^{\mathrm{proto}}}
\newcommand{\ProtoSection}[1]{\Gamma_{#1}^{\mathrm{proto}}}
\newcommand{\ProtoShell}[1]{\mathcal Z_{#1}^{\mathrm{proto}}}
\newcommand{\ProtoTransportSec}[1]{\mathfrak s_{#1}^{\mathrm{proto,sec}}}

\newcommand{\ProtoSectionPackage}{\mathfrak Q^{\mathrm{sub,proto,secgen}}}
\newcommand{\ProtoSectionState}[1]{\mathcal Q_{#1}^{\mathrm{psec}}}
\newcommand{\ProtoSectionBody}[1]{\mathcal B_{#1}^{\mathrm{psec}}}
\newcommand{\ProtoSectionProj}[1]{\Pi_{#1}^{\mathrm{psec}}}
\newcommand{\ProtoSectionFiber}[2]{\mathcal F_{#1}^{\mathrm{psec}}(#2)}
\newcommand{\ProtoSectionFunctional}[1]{\mathscr E_{#1}^{\mathrm{psec}}}
\newcommand{\ProtoSectionShell}[1]{\mathcal Z_{#1}^{\mathrm{psec}}}
\newcommand{\ProtoSectionMin}[1]{\mathfrak m_{#1}^{\mathrm{psec}}}
\newcommand{\InducedTransportSec}[1]{\widehat{\mathfrak s}_{#1}^{\mathrm{proto,sec}}}

\newcommand{\PreProtoSectionPackage}{\mathfrak R^{\mathrm{sub,proto,presec}}}
\newcommand{\PreProtoSectionState}[1]{\mathcal R_{#1}^{\mathrm{ppsec}}}
\newcommand{\PreProtoSectionBody}[1]{\mathcal B_{#1}^{\mathrm{ppsec}}}
\newcommand{\PreProtoSectionProj}[1]{\Pi_{#1}^{\mathrm{ppsec}}}
\newcommand{\PreProtoSectionFiber}[2]{\mathcal F_{#1}^{\mathrm{ppsec}}(#2)}
\newcommand{\PreProtoSectionProtoMin}[1]{\mathfrak n_{#1}^{\mathrm{ppsec}}}
\newcommand{\PreProtoSectionFunctional}[1]{\mathscr F_{#1}^{\mathrm{ppsec}}}
\newcommand{\PreProtoSectionShell}[1]{\mathcal Z_{#1}^{\mathrm{ppsec}}}

\newcommand{\PreProtoSectionImageCore}{\mathfrak R^{\mathrm{sub,proto,presec,img}}}
\newcommand{\PreProtoSectionImgBody}[1]{\mathcal B_{#1}^{\mathrm{ppsec,img}}}
\newcommand{\PreProtoSectionImgProj}[1]{\Pi_{#1}^{\mathrm{ppsec,img}}}
\newcommand{\PreProtoSectionImgFunctional}[1]{\mathscr J_{#1}^{\mathrm{ppsec,img}}}
\newcommand{\PreProtoSectionImgShell}[1]{\mathcal Z_{#1}^{\mathrm{ppsec,img}}}

\newtheorem{definition}{Definition}
\newtheorem{proposition}{Proposition}
\newtheorem{remark}{Remark}
\newtheorem{corollary}{Corollary}
\newtheorem{theorem}{Theorem}

\title{\textbf{The \TheoryName{}}\\[0.4em]
\large Primitive-Reservoir Observer-Reduction\\
\large Proto-Section-Generating Substrate Pre-Proto-Section Dynamics\\
\large Collapse to an Observer-Relevant\\
\large Pre-Proto-Section Minimizer-Image Core}
\author{}
\date{}

\begin{document}

\maketitle

\begin{abstract}
The current primitive-reservoir observer-reduction frontier already moved the live burden below the exact-shell-transport-section-generating substrate proto-section dynamics package \(\ProtoSectionPackage\): the active-primary theorem showed that the current proto-section package is canonically generated by a deeper proto-section-generating substrate pre-proto-section dynamics package \(\PreProtoSectionPackage\). Under the recorded safeguard against recursive same-output depth drift, however, the honest next move is no longer another deeper relay that merely preserves that same pre-proto-section package. The sharper benchmark-facing question is whether the full recorded pre-proto-section package still carries benchmark-neutral ambient freedom once the induced proto-section package is held fixed.

The present note proves that it does. It isolates the \emph{observer-relevant pre-proto-section minimizer-image core}: the minimizer-image compact body over the recorded proto-section body, the restricted projection back to that body, the restricted pre-proto-section functional on that image, the exact-shell image of the recorded proto-section shell, and benchmark faithfulness. The full exact pre-proto-section shell is then shown to be forced by that image core, and ambient pre-proto-section state-space directions outside the minimizer image do not add independent benchmark-facing freedom once the induced proto-section package is held fixed.

The benchmark boundary remains fixed throughout. The one operative metric is still exactly \(\CompletedMetric\), the ordinary matter route is unchanged, there is no second operative metric, the low-energy matter law is not enlarged, and no stronger promotion language is licensed. The positive result is therefore a genuine smaller theorem-level collapse strictly inside the current pre-proto-section frontier: the live burden contracts from the full pre-proto-section package \(\PreProtoSectionPackage\) to the smaller observer-relevant pre-proto-section minimizer-image core \(\PreProtoSectionImageCore\). What remains open is to derive or justify that smaller image core from still more primitive substrate structure, or else prove a still sharper collapse strictly inside it.
\end{abstract}

\section{Introduction}

The active derivational program of \TheoryName{} has again reached a sharper burden than another same-output deeper relay. The current active-primary theorem already removed the exact-shell-transport-section-generating substrate proto-section dynamics package \(\ProtoSectionPackage\) as the live primitive stopping point by showing that, once a benchmark-faithful pre-proto-section package \(\PreProtoSectionPackage\) is fixed, the recorded proto-section body, functional, and exact shell are its canonical fiberwise effective output. That theorem matters precisely because it blocks a familiar form of recursive depth drift: depth alone no longer counts once the live benchmark-relevant package datum is unchanged.

The next honest move is therefore not to relay the same pre-proto-section package through another deeper origin while preserving that package and its same downstream outputs. It is to ask whether the present pre-proto-section package itself still contains benchmark-neutral ambient freedom that can be removed without changing the induced proto-section package or the downstream chain depending on it. The present note takes that sharper internal-collapse route. It shows that the full ambient pre-proto-section state space and off-image body directions are not independently live once one retains only the minimizer image over the recorded proto-section body together with the restricted functional and exact-shell image it already determines.

\paragraph{Framework and claim boundary.}
\TheoryProgram{} denotes the broader \Aether{} / \Aflow{} framework built on the claims that the \Aether{} is the underlying four-dimensional substrate of reality, the \Aflow{} is its intrinsic ordered motion, observed three-dimensional space is the local experiential slice of that deeper substrate, S-time is the experienced order of change arising from matter, light, and the \Aflow{}, and observed expansion is the three-dimensional appearance of deeper four-dimensional ordered motion. Within that framework, \TheoryName{} denotes the exact-closure theory in which the effective gravitational dynamics are adopted to be exactly Einsteinian with universal matter coupling. In the active sequence, \emph{adoption} means use of that established relativistic dynamics without claiming substrate derivation, while \emph{derivation} is reserved for a first-principles recovery from explicit substrate variables.

The present note stays strictly inside that boundary. It does not introduce a second operative metric, it does not enlarge the low-energy matter law, and it does not reopen the already-generated proto-section package, the already-generated section core, or older seed, root, foundation, ground, origin, precursor, lift, shell-restriction, selector-law, combined-response, ambient-package, trace-core, or exact-shell-affine-nucleus layers as if they were again the live burden. Its narrower benchmark-facing role is to isolate the smaller current datum inside the pre-proto-section package itself.

\section{Fixed Inputs}

\begin{definition}[Recorded primitive continuation data]
\label{def:fixed}
Fix the following continuation data from the active primitive-reservoir observer-reduction line.
\begin{enumerate}
    \item For each sector label \(\sigma\in\{\Car,\Cur,\Hom\}\), a primitive forcing shell
    \[
    \ForSpace{\sigma},
    \]
    a visible reduction map
    \[
    \Reduction{\sigma}:\ForSpace{\sigma}\to X_{\sigma},
    \]
    and shell-level current and equilibrium carriers
    \[
    \CurrentCarrier{\sigma}:\ForSpace{\sigma}\to Z_{\sigma}^{\mathrm{cur}},
    \qquad
    \EqCarrier{\sigma}:\ForSpace{\sigma}\to Z_{\sigma}^{\mathrm{eq}}.
    \]
    \item The product shell
    \[
    \PrimShell
    \equiv
    \ForSpace{\Car}\times \ForSpace{\Cur}\times \ForSpace{\Hom}.
    \]
    \item The recorded metric and ordinary-matter outputs
    \[
    \MetricOut:\PrimShell\to\{\text{observer metrics}\},
    \qquad
    \MatterOut:\PrimShell\times\MatterSpace\to\{\text{observer stress tensors}\},
    \]
    satisfying
    \begin{equation}
    \MetricOut(\mathbf w)=\CompletedMetric,
    \qquad
    \MatterOut(\mathbf w,\psi)
    =
    T_{\mu\nu}^{\mathrm{matter}}[\CompletedMetric,\psi]
    \label{eq:fixedoutputs}
    \end{equation}
    for every \(\mathbf w\in\PrimShell\) and every matter configuration \(\psi\in\MatterSpace\).
\end{enumerate}
\end{definition}

\begin{definition}[Recorded downstream seed package]
\label{def:seed}
Fix \(\sigma\in\{\Car,\Cur,\Hom\}\). The active line already fixes the downstream seed package at current benchmark-facing scope:
\begin{enumerate}
    \item a finite-dimensional seed state space \(\SeedState{\sigma}\), a compact convex seed body \(\SeedBody{\sigma}\subset\SeedState{\sigma}\), the exact seed shell \(\SeedShell{\sigma}\subset\SeedBody{\sigma}\), and the exact seed support shell \(\SeedSupportShell{\sigma}\subset\SeedShell{\sigma}\);
    \item a seed zero core \(\SeedZeroCore{\sigma}\subset\SeedSupportShell{\sigma}\), the same hidden fiber space \(\HiddenSpace{\sigma}\) with compact hidden body \(\HiddenBody{\sigma}\subset\HiddenSpace{\sigma}\), and the fixed seed split
    \[
    \SeedSplit{\sigma}:
    \SeedZeroCore{\sigma}\times\HiddenBody{\sigma}
    \xrightarrow{\cong}
    \SeedSupportShell{\sigma};
    \]
    \item the seed hidden charge
    \[
    \SeedCharge{\sigma}:\HiddenSpace{\sigma}\to\EqnTarget{\sigma},
    \]
    a smooth observer-compatible exact seed zero-response realization
    \[
    \SeedEmbed{\sigma}:\ForSpace{\sigma}\hookrightarrow\SeedZeroCore{\sigma},
    \]
    and the fixed-data solvability / coercive package recorded through \(\SeedKernel{\sigma}\) and \(\SeedCoercive{\sigma}\);
    \item benchmark faithfulness, meaning preservation of the benchmark outputs \eqref{eq:fixedoutputs}, no second operative metric, no enlarged low-energy matter law, no change to the ordinary matter route, and no premature promotion from adoption to completed derivation.
\end{enumerate}
\end{definition}

\begin{definition}[Recorded current minimal germ zero-shell transport nucleus]
\label{def:nucleus}
Fix \(\sigma\in\{\Car,\Cur,\Hom\}\). The recorded current minimal germ zero-shell transport nucleus \(\GermNucleus\) for sector \(\sigma\) consists of:
\begin{enumerate}
    \item the germ state space \(\GermState{\sigma}\), compact convex germ body \(\GermBody{\sigma}\subset\GermState{\sigma}\), projection
    \[
    \GermProj{\sigma}:\GermState{\sigma}\to\SeedState{\sigma},
    \]
    and seed section
    \[
    \GermSeedSec{\sigma}:\SeedBody{\sigma}\to\GermBody{\sigma}
    \]
    satisfying
    \[
    \GermProj{\sigma}\bigl(\GermBody{\sigma}\bigr)=\SeedBody{\sigma},
    \qquad
    \GermProj{\sigma}\circ\GermSeedSec{\sigma}
    =
    \mathrm{id}_{\SeedBody{\sigma}};
    \]
    \item the restriction section
    \[
    \GermSection{\sigma}:\GermState{\sigma}\to\SubTarget{\sigma},
    \]
    and the exact germ shell
    \[
    \GermShell{\sigma}
    \equiv
    \bigl(\GermSection{\sigma}\bigr)^{-1}(0)\cap\GermBody{\sigma},
    \]
    with
    \[
    \GermProj{\sigma}\big|_{\GermShell{\sigma}}
    :
    \GermShell{\sigma}
    \xrightarrow{\cong}
    \SeedShell{\sigma};
    \]
    \item benchmark faithfulness, meaning preservation of the benchmark outputs \eqref{eq:fixedoutputs}, no second operative metric, no enlarged low-energy matter law, no change to the ordinary matter route, and no premature promotion from adoption to completed derivation.
\end{enumerate}
By the already-recorded current frontier theorem, the support shell, zero core, hidden split, hidden charge, exact zero-response realization, and fixed-data solvability / coercive package are then canonically induced downstream from this nucleus together with the fixed seed package of Definition \ref{def:seed}, rather than taken as additional independent benchmark-facing data.
\end{definition}

\begin{definition}[Recorded current germ-anchored exact-shell transport core]
\label{def:core}
Fix \(\sigma\in\{\Car,\Cur,\Hom\}\). The recorded current germ-anchored exact-shell transport core \(\ProtoTransportCore\) for sector \(\sigma\) consists of:
\begin{enumerate}
    \item the proto-germ state space \(\ProtoState{\sigma}\), compact convex proto-germ body \(\ProtoBody{\sigma}\subset\ProtoState{\sigma}\), and affine surjection
    \[
    \ProtoProj{\sigma}:\ProtoState{\sigma}\to\GermState{\sigma}
    \]
    satisfying
    \[
    \ProtoProj{\sigma}\bigl(\ProtoBody{\sigma}\bigr)=\GermBody{\sigma};
    \]
    \item the restriction section
    \[
    \ProtoSection{\sigma}:\ProtoState{\sigma}\to\SubTarget{\sigma},
    \]
    the exact proto-germ shell
    \[
    \ProtoShell{\sigma}
    \equiv
    \bigl(\ProtoSection{\sigma}\bigr)^{-1}(0)\cap\ProtoBody{\sigma},
    \]
    and the exact-shell transport diffeomorphism
    \[
    \ProtoProj{\sigma}\big|_{\ProtoShell{\sigma}}
    :
    \ProtoShell{\sigma}
    \xrightarrow{\cong}
    \GermShell{\sigma};
    \]
    \item benchmark faithfulness, meaning preservation of the benchmark outputs \eqref{eq:fixedoutputs}, no second operative metric, no enlarged low-energy matter law, no change to the ordinary matter route, and no premature promotion from adoption to completed derivation.
\end{enumerate}
\end{definition}

\begin{definition}[Recorded current germ-anchored exact-shell transport-section core]
\label{def:sectioncore}
Fix \(\sigma\in\{\Car,\Cur,\Hom\}\) and the recorded transport core of Definition \ref{def:core}. The recorded current germ-anchored exact-shell transport-section core \(\ProtoSectionCore\) for sector \(\sigma\) consists of:
\begin{enumerate}
    \item the canonical exact-shell transport section
    \[
    \ProtoTransportSec{\sigma}
    \equiv
    \left(
    \ProtoProj{\sigma}\big|_{\ProtoShell{\sigma}}
    \right)^{-1}
    :
    \GermShell{\sigma}\to\ProtoShell{\sigma}\subset\ProtoBody{\sigma}\subset\ProtoState{\sigma};
    \]
    \item benchmark faithfulness in the sense of Definition \ref{def:core}.
\end{enumerate}
The exact proto shell and the restricted shell projection back to \(\GermShell{\sigma}\) are understood as derived data rather than as separately primitive section-core data.
\end{definition}

\section{Recorded Frontier Package}

\begin{definition}[Recorded exact-shell-transport-section-generating substrate proto-section dynamics]
\label{def:protosection}
Fix \(\sigma\in\{\Car,\Cur,\Hom\}\). The recorded exact-shell-transport-section-generating substrate proto-section dynamics package \(\ProtoSectionPackage\) for sector \(\sigma\) consists of the following data.
\begin{enumerate}
    \item A finite-dimensional Euclidean proto-section state space \(\ProtoSectionState{\sigma}\), a compact convex proto-section body
    \[
    \ProtoSectionBody{\sigma}\subset\ProtoSectionState{\sigma}
    \]
    with nonempty interior, and an affine surjection
    \[
    \ProtoSectionProj{\sigma}:\ProtoSectionState{\sigma}\to\ProtoState{\sigma}
    \]
    satisfying
    \[
    \ProtoSectionProj{\sigma}\bigl(\ProtoSectionBody{\sigma}\bigr)=\ProtoBody{\sigma}.
    \]
    For each \(x\in\GermBody{\sigma}\), write
    \[
    \ProtoSectionFiber{\sigma}{x}
    \equiv
    \left\{
    y\in\ProtoSectionBody{\sigma}:
    \ProtoProj{\sigma}\bigl(\ProtoSectionProj{\sigma}(y)\bigr)=x
    \right\}.
    \]
    Each such fiber is required to be nonempty, compact, and convex.
    \item A nonnegative smooth fiber-selection functional
    \[
    \ProtoSectionFunctional{\sigma}:\ProtoSectionState{\sigma}\to[0,\infty)
    \]
    such that, for every \(x\in\GermBody{\sigma}\), the restriction of \(\ProtoSectionFunctional{\sigma}\) to \(\ProtoSectionFiber{\sigma}{x}\) is strictly convex and attains a unique minimizer. The resulting minimizer depends smoothly on \(x\).
    \item A closed embedded exact proto-section shell
    \[
    \ProtoSectionShell{\sigma}\subset\ProtoSectionBody{\sigma}
    \]
    satisfying
    \[
    \ProtoSectionShell{\sigma}
    =
    \left\{
    y\in\ProtoSectionBody{\sigma}:
    \ProtoProj{\sigma}\bigl(\ProtoSectionProj{\sigma}(y)\bigr)\in\GermShell{\sigma},
    \;
    \ProtoSectionFunctional{\sigma}(y)
    =
    \min_{z\in\ProtoSectionFiber{\sigma}{\ProtoProj{\sigma}(\ProtoSectionProj{\sigma}(y))}}
    \ProtoSectionFunctional{\sigma}(z)
    \right\},
    \]
    and
    \[
    \ProtoSectionProj{\sigma}\big|_{\ProtoSectionShell{\sigma}}
    :
    \ProtoSectionShell{\sigma}
    \xrightarrow{\cong}
    \ProtoShell{\sigma}
    \]
    is a diffeomorphism.
    \item Benchmark faithfulness, meaning preservation of the benchmark outputs \eqref{eq:fixedoutputs}, no second operative metric, no enlarged low-energy matter law, no change to the ordinary matter route, and no premature promotion from adoption to completed derivation.
\end{enumerate}
\end{definition}

\begin{definition}[Recorded proto-section-generating substrate pre-proto-section dynamics]
\label{def:preprotosection}
Fix \(\sigma\in\{\Car,\Cur,\Hom\}\). The recorded current pre-proto-section dynamics package \(\PreProtoSectionPackage\) for sector \(\sigma\) consists of the following data.
\begin{enumerate}
    \item A finite-dimensional Euclidean pre-proto-section state space \(\PreProtoSectionState{\sigma}\), a compact convex pre-proto-section body
    \[
    \PreProtoSectionBody{\sigma}\subset\PreProtoSectionState{\sigma}
    \]
    with nonempty interior, and an affine surjection
    \[
    \PreProtoSectionProj{\sigma}:\PreProtoSectionState{\sigma}\to\ProtoSectionState{\sigma}
    \]
    satisfying
    \[
    \PreProtoSectionProj{\sigma}\bigl(\PreProtoSectionBody{\sigma}\bigr)=\ProtoSectionBody{\sigma}.
    \]
    For each \(y\in\ProtoSectionBody{\sigma}\), write
    \[
    \PreProtoSectionFiber{\sigma}{y}
    \equiv
    \bigl(\PreProtoSectionProj{\sigma}\bigr)^{-1}(y)\cap\PreProtoSectionBody{\sigma}.
    \]
    Each such fiber is required to be nonempty, compact, and convex.
    \item A nonnegative smooth pre-proto-section functional
    \[
    \PreProtoSectionFunctional{\sigma}:\PreProtoSectionState{\sigma}\to[0,\infty)
    \]
    such that, for every \(y\in\ProtoSectionBody{\sigma}\), the restriction of \(\PreProtoSectionFunctional{\sigma}\) to \(\PreProtoSectionFiber{\sigma}{y}\) is strictly convex and attains a unique minimizer. The resulting minimizer depends smoothly on \(y\); write it as
    \[
    \PreProtoSectionProtoMin{\sigma}:\ProtoSectionBody{\sigma}\to\PreProtoSectionBody{\sigma}.
    \]
    These data satisfy
    \[
    \PreProtoSectionProj{\sigma}\circ\PreProtoSectionProtoMin{\sigma}
    =
    \mathrm{id}_{\ProtoSectionBody{\sigma}},
    \]
    and the effective minimum value recovers the recorded proto-section functional on the benchmark-relevant body:
    \[
    \ProtoSectionFunctional{\sigma}(y)
    =
    \PreProtoSectionFunctional{\sigma}\bigl(\PreProtoSectionProtoMin{\sigma}(y)\bigr)
    =
    \min_{u\in\PreProtoSectionFiber{\sigma}{y}}
    \PreProtoSectionFunctional{\sigma}(u)
    \qquad
    \text{for every }
    y\in\ProtoSectionBody{\sigma}.
    \]
    \item A closed embedded exact pre-proto-section shell
    \[
    \PreProtoSectionShell{\sigma}\subset\PreProtoSectionBody{\sigma}
    \]
    satisfying
    \[
    \PreProtoSectionShell{\sigma}
    =
    \left\{
    u\in\PreProtoSectionBody{\sigma}:
    \begin{array}{l}
    \PreProtoSectionProj{\sigma}(u)\in\ProtoSectionShell{\sigma},\\[0.2em]
    \PreProtoSectionFunctional{\sigma}(u)
    =
    \displaystyle\min_{v\in\PreProtoSectionFiber{\sigma}{\PreProtoSectionProj{\sigma}(u)}}
    \PreProtoSectionFunctional{\sigma}(v)
    \end{array}
    \right\},
    \]
    and
    \[
    \PreProtoSectionProj{\sigma}\big|_{\PreProtoSectionShell{\sigma}}
    :
    \PreProtoSectionShell{\sigma}
    \xrightarrow{\cong}
    \ProtoSectionShell{\sigma}
    \]
    is a diffeomorphism.
    \item Benchmark faithfulness, meaning preservation of the benchmark outputs \eqref{eq:fixedoutputs}, no second operative metric, no enlarged low-energy matter law, and presentation as a deeper substrate-side source for the current proto-section package rather than as a replacement benchmark theory.
\end{enumerate}
\end{definition}

\begin{remark}
Definitions \ref{def:protosection} and \ref{def:preprotosection} fix the current frontier package. The present note does not spend its move on another deeper derivation below that frontier. It asks instead whether the full recorded pre-proto-section package carries benchmark-neutral ambient freedom that can be removed without changing the induced proto-section package or the downstream chain that depends on it.
\end{remark}

\section{Observer-Relevant Pre-Proto-Section Minimizer-Image Core}

\begin{definition}[Observer-relevant pre-proto-section minimizer-image core]
\label{def:imagecore}
Fix \(\sigma\in\{\Car,\Cur,\Hom\}\) and a benchmark-faithful pre-proto-section package in the sense of Definition \ref{def:preprotosection}. The \emph{observer-relevant pre-proto-section minimizer-image core} \(\PreProtoSectionImageCore\) for sector \(\sigma\) consists of:
\begin{enumerate}
    \item the compact minimizer-image body
    \[
    \PreProtoSectionImgBody{\sigma}
    \equiv
    \PreProtoSectionProtoMin{\sigma}\bigl(\ProtoSectionBody{\sigma}\bigr)
    \subset
    \PreProtoSectionBody{\sigma},
    \]
    together with the restricted projection
    \[
    \PreProtoSectionImgProj{\sigma}
    \equiv
    \PreProtoSectionProj{\sigma}\big|_{\PreProtoSectionImgBody{\sigma}}
    :
    \PreProtoSectionImgBody{\sigma}\to\ProtoSectionBody{\sigma};
    \]
    \item the restricted functional
    \[
    \PreProtoSectionImgFunctional{\sigma}
    \equiv
    \PreProtoSectionFunctional{\sigma}\big|_{\PreProtoSectionImgBody{\sigma}};
    \]
    \item the exact-shell image
    \[
    \PreProtoSectionImgShell{\sigma}
    \equiv
    \PreProtoSectionProtoMin{\sigma}\bigl(\ProtoSectionShell{\sigma}\bigr)
    \subset
    \PreProtoSectionShell{\sigma};
    \]
    \item benchmark faithfulness in the sense of Definition \ref{def:preprotosection}(4).
\end{enumerate}
\end{definition}

\begin{remark}
Definition \ref{def:imagecore} is intentionally smaller than Definition \ref{def:preprotosection}. It keeps only the minimizer image over the recorded proto-section body together with the restricted functional and the exact-shell image it already determines. The ambient pre-proto-section state space and off-image body directions are not taken as independent core data.
\end{remark}

\section{Collapse of the Full Pre-Proto-Section Package}

\begin{proposition}[The minimizer-image body carries the body-level content]
\label{prop:body}
For every benchmark-faithful pre-proto-section package, \(\PreProtoSectionImgBody{\sigma}\) is compact and the restricted projection
\[
\PreProtoSectionImgProj{\sigma}:\PreProtoSectionImgBody{\sigma}\to\ProtoSectionBody{\sigma}
\]
is a homeomorphism with inverse \(\PreProtoSectionProtoMin{\sigma}\).
\end{proposition}

\begin{proof}
The body \(\ProtoSectionBody{\sigma}\) is compact by Definition \ref{def:protosection}, and \(\PreProtoSectionProtoMin{\sigma}\) is continuous by Definition \ref{def:preprotosection}(2). Therefore
\[
\PreProtoSectionImgBody{\sigma}
=
\PreProtoSectionProtoMin{\sigma}\bigl(\ProtoSectionBody{\sigma}\bigr)
\]
is compact. The identity
\[
\PreProtoSectionProj{\sigma}\circ\PreProtoSectionProtoMin{\sigma}
=
\mathrm{id}_{\ProtoSectionBody{\sigma}}
\]
shows that \(\PreProtoSectionImgProj{\sigma}\) is surjective and that \(\PreProtoSectionProtoMin{\sigma}\) is a right inverse for it. If \(u_{1},u_{2}\in\PreProtoSectionImgBody{\sigma}\) satisfy
\[
\PreProtoSectionImgProj{\sigma}(u_{1})=\PreProtoSectionImgProj{\sigma}(u_{2})=y,
\]
then each \(u_{i}\) lies in the fiber \(\PreProtoSectionFiber{\sigma}{y}\) and is the unique minimizer there, so \(u_{1}=u_{2}=\PreProtoSectionProtoMin{\sigma}(y)\). Thus \(\PreProtoSectionImgProj{\sigma}\) is bijective with inverse \(\PreProtoSectionProtoMin{\sigma}\). Since \(\PreProtoSectionImgBody{\sigma}\) is compact and \(\ProtoSectionBody{\sigma}\) is Hausdorff, the continuous bijection \(\PreProtoSectionImgProj{\sigma}\) is a homeomorphism.
\end{proof}

\begin{proposition}[The recorded proto-section functional is carried by the image core]
\label{prop:functional}
For every benchmark-faithful pre-proto-section package, the recorded proto-section functional is recovered entirely from the restricted image-core functional:
\[
\ProtoSectionFunctional{\sigma}(y)
=
\PreProtoSectionImgFunctional{\sigma}\bigl(\PreProtoSectionProtoMin{\sigma}(y)\bigr)
\qquad
\text{for every }y\in\ProtoSectionBody{\sigma}.
\]
\end{proposition}

\begin{proof}
By Definition \ref{def:imagecore}(2),
\[
\PreProtoSectionImgFunctional{\sigma}\bigl(\PreProtoSectionProtoMin{\sigma}(y)\bigr)
=
\PreProtoSectionFunctional{\sigma}\bigl(\PreProtoSectionProtoMin{\sigma}(y)\bigr).
\]
Definition \ref{def:preprotosection}(2) then gives
\[
\PreProtoSectionFunctional{\sigma}\bigl(\PreProtoSectionProtoMin{\sigma}(y)\bigr)
=
\ProtoSectionFunctional{\sigma}(y)
\]
for every \(y\in\ProtoSectionBody{\sigma}\).
\end{proof}

\begin{proposition}[The full exact pre-proto-section shell is forced by the image core]
\label{prop:shell}
For every benchmark-faithful pre-proto-section package,
\[
\PreProtoSectionShell{\sigma}
=
\PreProtoSectionImgShell{\sigma}
=
\PreProtoSectionProtoMin{\sigma}\bigl(\ProtoSectionShell{\sigma}\bigr),
\]
and the restricted projection
\[
\PreProtoSectionImgProj{\sigma}\big|_{\PreProtoSectionImgShell{\sigma}}
:
\PreProtoSectionImgShell{\sigma}\xrightarrow{\cong}\ProtoSectionShell{\sigma}
\]
is a diffeomorphism.
\end{proposition}

\begin{proof}
If \(u\in\PreProtoSectionShell{\sigma}\), then Definition \ref{def:preprotosection}(3) gives
\[
y\equiv \PreProtoSectionProj{\sigma}(u)\in\ProtoSectionShell{\sigma},
\]
and \(u\) minimizes \(\PreProtoSectionFunctional{\sigma}\) on the fiber \(\PreProtoSectionFiber{\sigma}{y}\). By Definition \ref{def:preprotosection}(2), that minimizer is unique and equals \(\PreProtoSectionProtoMin{\sigma}(y)\). Hence
\[
u=\PreProtoSectionProtoMin{\sigma}(y)\in\PreProtoSectionImgShell{\sigma},
\]
so \(\PreProtoSectionShell{\sigma}\subset\PreProtoSectionImgShell{\sigma}\).

Conversely, if \(y\in\ProtoSectionShell{\sigma}\), then \(\PreProtoSectionProtoMin{\sigma}(y)\) is by definition the unique minimizer of \(\PreProtoSectionFunctional{\sigma}\) on the fiber over \(y\). Therefore \(\PreProtoSectionProtoMin{\sigma}(y)\in\PreProtoSectionShell{\sigma}\), so \(\PreProtoSectionImgShell{\sigma}\subset\PreProtoSectionShell{\sigma}\). This proves the equality of shells.

Finally,
\[
\PreProtoSectionImgProj{\sigma}\big|_{\PreProtoSectionImgShell{\sigma}}
=
\PreProtoSectionProj{\sigma}\big|_{\PreProtoSectionShell{\sigma}},
\]
and the latter is a diffeomorphism onto \(\ProtoSectionShell{\sigma}\) by Definition \ref{def:preprotosection}(3).
\end{proof}

\begin{proposition}[The image core canonically reproduces the current proto-section package]
\label{prop:package}
Fix a benchmark-faithful pre-proto-section package. Then the current proto-section package of Definition \ref{def:protosection} is exactly the benchmark-relevant package induced by the smaller image core. In particular:
\begin{enumerate}
    \item the recorded proto-section body and its body-level transport over that body are recovered by Proposition \ref{prop:body};
    \item the recorded proto-section functional on \(\ProtoSectionBody{\sigma}\) is recovered by Proposition \ref{prop:functional};
    \item the recorded exact proto-section shell is recovered by Proposition \ref{prop:shell}.
\end{enumerate}
\end{proposition}

\begin{proof}
Item (1) is Proposition \ref{prop:body}. Item (2) is Proposition \ref{prop:functional}. Item (3) is Proposition \ref{prop:shell}. These statements are exactly the benchmark-relevant constituents of the current proto-section package at current scope.
\end{proof}

\begin{proposition}[All downstream recovery factors through the image core]
\label{prop:downstream}
Fix a benchmark-faithful pre-proto-section package. Then the current germ-anchored exact-shell transport-section core \(\ProtoSectionCore\) and the previously recorded downstream seed, root, foundation, ground, origin, precursor, lift, shell-restriction, split-core, selector-law, combined-response, ambient-package, trace-core, exact-shell-affine-nucleus, and one-metric observer-package consequences on the active line all factor through \(\PreProtoSectionImageCore\).
\end{proposition}

\begin{proof}
The previously recorded theorem deriving the current germ-anchored exact-shell transport-section core \(\ProtoSectionCore\) requires exactly the proto-section package of Definition \ref{def:protosection} as its live input at current scope. Proposition \ref{prop:package} shows that this entire proto-section package is already determined by the smaller image core. Therefore the recorded section-core recovery and the previously recorded downstream chain that depends on that proto-section package do not require the full ambient pre-proto-section package as additional independent data; they factor through \(\PreProtoSectionImageCore\).
\end{proof}

\begin{proposition}[Ambient pre-proto-section directions are benchmark neutral at current scope]
\label{prop:neutral}
For every benchmark-faithful pre-proto-section package, the ambient pre-proto-section state space \(\PreProtoSectionState{\sigma}\) and the off-image body directions
\[
\PreProtoSectionBody{\sigma}\setminus\PreProtoSectionImgBody{\sigma}
\]
do not add independent benchmark-facing freedom beyond \(\PreProtoSectionImageCore\).
\end{proposition}

\begin{proof}
By Proposition \ref{prop:body}, the full body-level content relevant to the recorded proto-section package is already fixed by the minimizer-image body \(\PreProtoSectionImgBody{\sigma}\) and the restricted projection \(\PreProtoSectionImgProj{\sigma}\). By Proposition \ref{prop:functional}, the current proto-section functional is already fixed by the restricted image-core functional \(\PreProtoSectionImgFunctional{\sigma}\). By Proposition \ref{prop:shell}, the full exact pre-proto-section shell is not additional data beyond the shell image; it is exactly that image. By Proposition \ref{prop:downstream}, all downstream recovery of the current section core and the previously recorded active-line consequences factors through this smaller core. Finally, benchmark faithfulness is already part of Definition \ref{def:imagecore}(4) and preserves the fixed outputs \eqref{eq:fixedoutputs}. Therefore nothing benchmark-facing at the present frontier requires the ambient pre-proto-section state-space presentation or body points outside the minimizer image as additional independent data.
\end{proof}

\begin{proposition}[The benchmark package remains fixed]
\label{prop:benchmark}
Every observer-relevant pre-proto-section minimizer-image core preserves the same benchmark one-metric package \eqref{eq:fixedoutputs}. In particular, the operative metric remains exactly \(\CompletedMetric\), the ordinary matter route is unchanged, there is no second operative metric, and the low-energy matter law is not enlarged.
\end{proposition}

\begin{proof}
This is part of benchmark faithfulness in Definition \ref{def:imagecore}(4). The present note only removes benchmark-neutral ambient pre-proto-section freedom from the live burden. It does not alter the benchmark exact-closure package.
\end{proof}

\section{Main Theorem}

\begin{theorem}[Proto-section-generating substrate pre-proto-section dynamics collapse to an observer-relevant pre-proto-section minimizer-image core]
\label{thm:main}
Fix the primitive continuation data of Definition \ref{def:fixed}, the recorded downstream seed package of Definition \ref{def:seed}, the recorded current germ nucleus of Definition \ref{def:nucleus}, the recorded current transport core of Definition \ref{def:core}, the recorded current section core of Definition \ref{def:sectioncore}, the recorded proto-section package of Definition \ref{def:protosection}, and a benchmark-faithful pre-proto-section package in the sense of Definition \ref{def:preprotosection}. Then, for each sector \(\sigma\in\{\Car,\Cur,\Hom\}\), the live benchmark-facing burden inside the current pre-proto-section frontier collapses from the full pre-proto-section package \(\PreProtoSectionPackage\) to the smaller observer-relevant pre-proto-section minimizer-image core \(\PreProtoSectionImageCore\). More precisely:
\begin{enumerate}
    \item the body-level content of the full pre-proto-section package is already carried by the minimizer-image body \(\PreProtoSectionImgBody{\sigma}\), the restricted projection \(\PreProtoSectionImgProj{\sigma}\), and the restricted image-core functional \(\PreProtoSectionImgFunctional{\sigma}\) by Propositions \ref{prop:body} and \ref{prop:functional};
    \item the full exact pre-proto-section shell is forced to equal the minimizer image of the recorded proto-section shell by Proposition \ref{prop:shell};
    \item the current proto-section package and therefore the current germ-anchored exact-shell transport-section core \(\ProtoSectionCore\) together with the previously recorded downstream consequences all factor through \(\PreProtoSectionImageCore\) by Propositions \ref{prop:package} and \ref{prop:downstream};
    \item ambient pre-proto-section directions outside that image core are benchmark neutral at current scope by Proposition \ref{prop:neutral};
    \item the benchmark boundary remains fixed by Proposition \ref{prop:benchmark};
    \item consequently the remaining honest burden below the previous pre-proto-section frontier is no longer the full proto-section-generating substrate pre-proto-section dynamics package \(\PreProtoSectionPackage\) as such. What remains is to derive or justify the sharper image core \(\PreProtoSectionImageCore\) from still more primitive substrate structure, or else prove a still smaller theorem-level collapse strictly inside that image core.
\end{enumerate}
\end{theorem}

\begin{proof}
Item (1) is Propositions \ref{prop:body} and \ref{prop:functional}. Item (2) is Proposition \ref{prop:shell}. Item (3) is Propositions \ref{prop:package} and \ref{prop:downstream}. Item (4) is Proposition \ref{prop:neutral}. Item (5) is Proposition \ref{prop:benchmark}. Item (6) is the routing consequence of the previous items together with the active safeguard against counting another same-output deeper-origin relay as new front-line progress when the live benchmark-relevant datum is unchanged.
\end{proof}

\begin{corollary}[What must be done next]
\label{cor:next}
After Theorem \ref{thm:main}, the next honest benchmark-facing manuscript must either:
\begin{enumerate}
    \item derive or justify the observer-relevant pre-proto-section minimizer-image core \(\PreProtoSectionImageCore\) from still more primitive substrate structure in a way that changes benchmark-facing status;
    \item identify a genuinely smaller theorem-level observer-relevant datum strictly inside \(\PreProtoSectionImageCore\) and prove a sharper collapse to it;
    \item do either of the previous two while preserving the same benchmark one-metric package, the same ordinary matter route, and the same claim boundary.
\end{enumerate}
Another same-output deeper-origin relay that preserves this image core and therefore merely reproduces the same current proto-section package \(\ProtoSectionPackage\), the same current germ-anchored exact-shell transport-section core \(\ProtoSectionCore\), and the same previously recorded downstream consequences, another re-expansion back to the full pre-proto-section package as if its screened ambient directions were still live, another reopening of the already-screened proto-section package or older larger active-line presentations as if they were again the current burden, or any move that introduces a second operative metric, enlarges the ordinary matter law, or uses stronger promotion language does not satisfy the present burden.
\end{corollary}

\begin{proof}
Theorem \ref{thm:main} shows that the full pre-proto-section package is no longer the minimal benchmark-facing datum at this stage. Therefore a subsequent manuscript must improve on the smaller image core rather than merely restating the larger package.
\end{proof}

\section{Conclusion}

The positive result of this note is a disciplined compression step inside the current pre-proto-section frontier. The previous active-primary theorem already removed the current proto-section package as the primitive stopping point by deriving it from a deeper pre-proto-section package. The present theorem then takes the remaining honest internal route available there: it shows that the full pre-proto-section package is itself not the minimal benchmark-facing datum.

What survives as the live burden is the observer-relevant pre-proto-section minimizer-image core. That sharper datum consists of the minimizer-image compact body over the recorded proto-section body, the restricted projection back to that body, the restricted pre-proto-section functional on that image, the exact-shell image of the recorded proto-section shell, and benchmark faithfulness. The full exact pre-proto-section shell is forced to equal that image shell, and ambient pre-proto-section directions outside the minimizer image do not add independent benchmark-facing freedom once the induced proto-section package is fixed.

This is the honest next burden after the present theorem. A new primary manuscript must derive or justify that sharper image core from still more primitive substrate structure, or else collapse it further, while preserving the same benchmark one-metric package, the same ordinary matter route, and the same claim boundary. The current result does not yet complete a first-principles derivation of GR from substrate microphysics, but it does narrow the live derivational burden one level further on the active line.

\end{document}
